% ============================================================================
%  3.9 PRUEBAS
% ============================================================================

\section{Pruebas}

\subsection{Pruebas realizadas}

Se realizaron las siguientes pruebas para validar el funcionamiento
del m\'odulo:

\begin{enumerate}
  \item \textbf{Prueba 1: Acceso al m\'odulo.} Se verific\'o que
    los eventos se muestran en el detalle del escenario con el bot\'on
    de inscripci\'on visible.

  \item \textbf{Prueba 2: Inscripci\'on exitosa.} Un usuario
    autenticado toca ``Inscribirse'' en un evento. El bot\'on cambia
    a ``Inscrito'' con fondo azul. El conteo de inscritos se
    incrementa en 1.

  \item \textbf{Prueba 3: Inscripci\'on duplicada.} Se intent\'o
    inscribir dos veces al mismo evento. Se mostr\'o el error
    ``Ya est\'as inscrito en este evento''.

  \item \textbf{Prueba 4: Cancelar inscripci\'on.} Un usuario
    inscrito toca ``Inscrito''. Se muestra di\'alogo de
    confirmaci\'on. Al confirmar, el bot\'on vuelve a
    ``Inscribirse'' y el conteo disminuye.

  \item \textbf{Prueba 5: Ver mis inscripciones.} En la pantalla
    de perfil, la secci\'on ``Mis inscripciones'' muestra todas las
    inscripciones activas del usuario.

  \item \textbf{Prueba 6: Sin inscripciones.} Un usuario nuevo
    ve el mensaje ``No tienes inscripciones a\'un'' en el perfil.

  \item \textbf{Prueba 7: Persistencia.} Se inscribió, se cerr\'o
    la app, se reabrió. La inscripci\'on persiste correctamente.

  \item \textbf{Prueba 8: Usuario no autenticado.} Un usuario sin
    sesi\'on ve la alerta ``Debes iniciar sesi\'on para
    inscribirte''.

  \item \textbf{Prueba 9: Manejo de errores de red.} Se simul\'o
    una conexi\'on lenta. El bot\'on mostr\'o el indicador de
    progreso y al fallar mostr\'o la alerta de error.

  \item \textbf{Prueba 10: Verificaci\'on en Supabase.} Se
    verific\'o desde Supabase Studio que los registros se crean
    correctamente en la tabla \texttt{event\_registrations}.
\end{enumerate}

\subsection{Resultados de las pruebas}

\begin{table}[H]
\centering
\begin{tabular}{@{}cll@{}}
\toprule
\textbf{\#} & \textbf{Prueba} & \textbf{Resultado} \\
\midrule
1 & Acceso al m\'odulo & \checkmark\ Pass \\
2 & Inscripci\'on exitosa & \checkmark\ Pass \\
3 & Inscripci\'on duplicada & \checkmark\ Pass \\
4 & Cancelar inscripci\'on & \checkmark\ Pass \\
5 & Ver mis inscripciones & \checkmark\ Pass \\
6 & Sin inscripciones & \checkmark\ Pass \\
7 & Persistencia & \checkmark\ Pass \\
8 & Usuario no autenticado & \checkmark\ Pass \\
9 & Manejo de errores & \checkmark\ Pass \\
10 & Verificaci\'on en Supabase & \checkmark\ Pass \\
\bottomrule
\end{tabular}
\caption{Resultados de las pruebas del m\'odulo}
\end{table}

\subsection{TypeScript}

Se ejecut\'o \texttt{tsc --noEmit} y el proyecto compila sin
errores de tipo, lo que confirma la correcta integraci\'on de
los tipos \texttt{EventRegistration} y
\texttt{EventRegistrationCount} en toda la aplicaci\'on.
